\documentclass[11pt,a4paper]{scrartcl}

\usepackage{amssymb}
\usepackage{amsmath}

\usepackage[utf8]{inputenc}
\usepackage[T1]{fontenc}
\newcommand{\ats}[1]{\par\noindent\textit{Atsakymas:} #1}
\usepackage{cancel}
\usepackage{tikz}
\usepackage{lmodern} 
\usepackage{amsmath,amssymb,amsthm}
\usepackage{graphicx}
\usepackage[dvipsnames,svgnames]{xcolor}
\usepackage{hyperref}
\usepackage{mathrsfs}
\usepackage[shortlabels]{enumitem}
\usepackage[obeyFinal,textsize=scriptsize,shadow,loadshadowlibrary]{todonotes}
\usepackage{mathtools}
\usepackage{microtype}

%%% Paragraph layout
\setlength{\parindent}{0pt}
\setlength{\parskip}{0.6\baselineskip}

%%% Colored sections
\renewcommand*{\sectionformat}%
  {\color{purple}\S\thesection\enskip}
\renewcommand*{\subsectionformat}%
  {\color{purple}\S\thesubsection\enskip}
\renewcommand*{\subsubsectionformat}%
  {\color{purple}\S\thesubsubsection\enskip}
\addtokomafont{paragraph}{\color{orange!35!black}\P\ }
\KOMAoptions{numbers=noenddot}

%%% Title
\addtokomafont{subtitle}{\Large}
\setkomafont{author}{\Large\scshape}
\setkomafont{date}{\Large\normalsize}
% Neigiamas vertikalus tarpas, kuris pakelia dokumento antraštę į viršų. 
% Galite keisti -2.5cm į -3cm (aukščiau) arba -1.5cm (žemiau).
\titlehead{\vspace*{-7.5cm}} 

% PRIDĖTA: Automatiškai perrašome datą į mokyklos pavadinimą
\AtBeginDocument{%
  \date{\normalsize Vilniaus licėjus}
}

%%% Page setup
\usepackage[headsepline]{scrlayer-scrpage}
\renewcommand{\headfont}{}
\addtolength{\textheight}{3.14cm}
\setlength{\footskip}{0.5in}
\setlength{\headsep}{10pt}
% Puslapio antraštė turi rodyti tik konkrečius dokumento duomenis.
% Nustatome juos aiškiai, kad neatsirastų "\@author" / "\@title" arba senų reikšmių.
\makeatletter
\AtBeginDocument{%
  \ihead{\footnotesize\textbf{Andrius Berniukevičius}}%
  \ohead{\footnotesize\textbf{\@title}}%
}
\makeatother
\automark{section}
\chead{}
\cfoot{\pagemark}

%%% Macros
\providecommand{\ol}{\overline}
\providecommand{\ul}{\underline}
\providecommand{\wt}{\widetilde}
\providecommand{\wh}{\widehat}
\providecommand{\eps}{\varepsilon}
\providecommand{\half}{\frac{1}{2}}
\providecommand{\inv}{^{-1}}
\newcommand{\dang}{\measuredangle} %% Directed angle
\newcommand{\vocab}[1]{\textbf{\color{blue}\sffamily #1}}
\providecommand{\CC}{\mathbb C}
\providecommand{\FF}{\mathbb F}
\providecommand{\NN}{\mathbb N}
\providecommand{\QQ}{\mathbb Q}
\providecommand{\RR}{\mathbb R}
\providecommand{\ZZ}{\mathbb Z}
\providecommand{\ts}{\textsuperscript}
\providecommand{\dg}{^\circ}
\providecommand{\ii}{\item}
\providecommand{\defeq}{\coloneqq}
\DeclareMathOperator*{\lcm}{lcm}
\DeclareMathOperator*{\argmin}{arg min}
\DeclareMathOperator*{\argmax}{arg max}
\providecommand{\hrulebar}{\par
\hspace{\fill}\rule{0.95\linewidth}{.7pt}\hspace{\fill}
\par\nointerlineskip \vspace{\baselineskip}}

%%% Optional colored theorems
\usepackage{thmtools}
\usepackage[framemethod=TikZ]{mdframed}
\mdfdefinestyle{mdbluebox}{%
  roundcorner=10pt,
  linewidth=1pt,
  skipabove=12pt,
  innerbottommargin=9pt,
  skipbelow=2pt,
  linecolor=blue,
  nobreak=true,
  backgroundcolor=TealBlue!5,
}
\declaretheoremstyle[
  headfont=\sffamily\bfseries\color{MidnightBlue},
  mdframed={style=mdbluebox},
  headpunct={\\[3pt]},
  postheadspace={0pt}
]{thmbluebox}
\mdfdefinestyle{mdredbox}{%
  linewidth=0.5pt,
  skipabove=12pt,
  frametitleaboveskip=5pt,
  frametitlebelowskip=0pt,
  skipbelow=2pt,
  frametitlefont=\bfseries,
  innertopmargin=4pt,
  innerbottommargin=8pt,
  nobreak=true,
  backgroundcolor=Salmon!5,
  linecolor=RawSienna,
}
\declaretheoremstyle[
  headfont=\bfseries\color{RawSienna},
  mdframed={style=mdredbox},
  headpunct={\\[3pt]},
  postheadspace={0pt},
]{thmredbox}
\mdfdefinestyle{mdgreenbox}{%
  skipabove=8pt,
  linewidth=2pt,
  rightline=false,
  leftline=true,
  topline=false,
  bottomline=false,
  linecolor=ForestGreen,
  backgroundcolor=ForestGreen!5,
}
\declaretheoremstyle[
  headfont=\bfseries\sffamily\color{ForestGreen!70!black},
  bodyfont=\normalfont,
  spaceabove=2pt,
  spacebelow=1pt,
  mdframed={style=mdgreenbox},
  headpunct={ --- },
]{thmgreenbox}
\mdfdefinestyle{mdblackbox}{%
  skipabove=8pt,
  linewidth=3pt,
  rightline=false,
  leftline=true,
  topline=false,
  bottomline=false,
  linecolor=black,
  backgroundcolor=RedViolet!5!gray!5,
}
\declaretheoremstyle[
  headfont=\bfseries,
  bodyfont=\normalfont\small,
  spaceabove=0pt,
  spacebelow=0pt,
  mdframed={style=mdblackbox}
]{thmblackbox}

%% Aplinkos su rėmeliais (thmtools)
\declaretheorem[style=thmbluebox,name=Teorema]{theorem}
\declaretheorem[style=thmbluebox,name=Lema]{lemma}
\declaretheorem[style=thmbluebox,name=Savybė]{property}
\declaretheorem[style=thmbluebox,name=Teiginys]{proposition}
\declaretheorem[style=thmbluebox,name=Išvada]{corollary}
\declaretheorem[style=thmbluebox,name=Teorema,numbered=no]{theorem*}
\declaretheorem[style=thmbluebox,name=Lema,numbered=no]{lemma*}
\declaretheorem[style=thmbluebox,name=Teiginys,numbered=no]{proposition*}
\declaretheorem[style=thmbluebox,name=Išvada,numbered=no]{corollary*}
\declaretheorem[style=thmgreenbox,name=Algoritmas]{algorithm}
\declaretheorem[style=thmgreenbox,name=Algoritmas,numbered=no]{algorithm*}
\declaretheorem[style=thmgreenbox,name=Teiginys]{claim}
\declaretheorem[style=thmgreenbox,name=Teiginys,numbered=no]{claim*}
\declaretheorem[style=thmredbox,name=Pavyzdys]{example}
\declaretheorem[style=thmredbox,name=Pavyzdys,numbered=no]{example*}
\declaretheorem[style=thmblackbox,name=Pastaba]{remark}
\declaretheorem[style=thmblackbox,name=Pastaba,numbered=no]{remark*}

%% Standartinės amsthm aplinkos (Apibrėžimai ir užduotys)
\theoremstyle{definition}
\ifcsname conjecture\endcsname\else\newtheorem{conjecture}{Hipotezė}\fi
\ifcsname definition\endcsname\else\newtheorem{definition}{Apibrėžimas}\fi
\ifcsname fact\endcsname\else\newtheorem{fact}{Faktas}\fi
\ifcsname answer\endcsname\else\newtheorem{answer}{Atsakymas}\fi
\ifcsname ques\endcsname\else\newtheorem{ques}{Klausimas}\fi
\ifcsname exercise\endcsname\else\newtheorem{exercise}{Užduotis}\fi
\ifcsname problem\endcsname\else\newtheorem{problem}{Uždavinys}[section]\fi
\ifcsname conjecture*\endcsname\else\newtheorem*{conjecture*}{Hipotezė}\fi
\ifcsname definition*\endcsname\else\newtheorem*{definition*}{Apibrėžimas}\fi
\ifcsname fact*\endcsname\else\newtheorem*{fact*}{Faktas}\fi
\ifcsname answer*\endcsname\else\newtheorem*{answer*}{Atsakymas}\fi
\ifcsname ques*\endcsname\else\newtheorem*{ques*}{Klausimas}\fi
\ifcsname exercise*\endcsname\else\newtheorem*{exercise*}{Užduotis}\fi
\ifcsname problem*\endcsname\else\newtheorem*{problem*}{Uždavinys}\fi

\newcounter{answerssection}
\newcommand{\answerssection}{%
  \setcounter{answerssection}{\value{section}}%
  \section{Atsakymai}%
  \setcounter{problem}{0}%
  \renewcommand{\theproblem}{\arabic{answerssection}.\arabic{problem}}%
}

%% GALUTINIS SPRENDIMAS PROOF APLINKAI
%% --- PRADŽIA: Įrodymo aplinkos sutvarkymas ---
\makeatletter

% 1. Užtikriname, kad žodis būtų "Įrodymas", net jei "babel" paketas bando jį perrašyti
\AtBeginDocument{%
  \renewcommand{\proofname}{Įrodymas}%
  \@ifpackageloaded{babel}{%
    \addto\captionslithuanian{\renewcommand{\proofname}{Įrodymas}}%
  }{}%
}

% 2. Priverstinai pašaliname scrartcl klasės bold šriftą ir paliekame TIK italic
% 3. Sujungiame su tuščiu kvadratėliu dešiniajame kampe.
\renewcommand{\qedsymbol}{\ensuremath{\Box}}
\renewenvironment{proof}[1][\proofname]{\par
  \pushQED{\qed}%
  \normalfont \topsep6\p@\@plus6\p@\relax
  \trivlist
  \item[\hskip\labelsep
        \normalfont\itshape % Priverčia naudoti tik pasvirą šriftą, kaip nuotraukoje
    #1\@addpunct{.}]\ignorespaces
}{%
  \popQED\endtrivlist\@endpefalse
}
\newenvironment{solution}{\par
  \normalfont \topsep6\p@\@plus6\p@\relax
  \trivlist
  \item[\hskip\labelsep
        \normalfont\itshape
    \textit{Sprendimas}\@addpunct{.}]\ignorespaces
}{%
  \endtrivlist\@endpefalse
}
\newcommand{\ans}[1]{\par\noindent\textit{Atsakymas:} #1}
\makeatother


\title{Skaičių teorija: Skaičių „DNR“ ir pirminiai skaičiai}
\author{Andrius Berniukevičius}

\newenvironment{solution}
    {\begin{list}{}{\setlength{\leftmargin}{10pt}\setlength{\rightmargin}{10pt}}\item[]}
    {\end{list}}

\begin{document}

\maketitle

\section{Skaičių „DNR“: Pagrindinė aritmetikos teorema}

Iki šiol mes mokėmės skaičius išardyti sudėties būdu ($100a + 10b + c$). Tačiau tikroji skaičiaus prigimtis slypi ne sudėtyje, o \textbf{daugyboje}.

Įsivaizduokite, kad visi sveikieji skaičiai yra pastatai. Jei skaičių galime padalinti, vadinasi, pastatą galime išardyti į mažesnius blokus. Pavyzdžiui, $60 = 6 \cdot 10$. Bet $6$ ir $10$ dar nėra mažiausios detalės. Ardome toliau: $6 = 2 \cdot 3$, o $10 = 2 \cdot 5$. 
Dabar turime $60 = 2 \cdot 2 \cdot 3 \cdot 5 = 2^2 \cdot 3 \cdot 5$. 

Skaičiai $2, 3$ ir $5$ toliau nebesiskaido. Tai yra \textbf{pirminiai skaičiai} – patys elementariausi matematikos „atomai“, iš kurių daugybos būdu pastatomas bet koks kitas skaičius.

\begin{center}
\fbox{
\begin{minipage}{0.9\textwidth}
\textbf{Pagrindinė aritmetikos teorema}\\
Kiekvienas natūralusis skaičius (didesnis už $1$) gali būti išskaidytas pirminių skaičių sandauga, ir toks skaidymas yra \textbf{vienintelis} (neskaitant dauginamųjų tvarkos). \\
Šis skaidymas $n = p_1^{a_1} \cdot p_2^{a_2} \cdot \dots \cdot p_k^{a_k}$ yra unikalus skaičiaus \textbf{„DNR kodas“}.
\end{minipage}
}
\end{center}

Kodėl tai galinga? Nes skaičiaus DNR pasako viską apie jo savybes. Pavyzdžiui, jei matome skaičių $x^2$ (tobuląjį kvadratą), jo DNR \textbf{visi} pirminių skaičių laipsniai privalo būti lyginiai ($x^2 = p_1^{2a_1} \cdot p_2^{2a_2} \dots$).

% ==========================================
% DALIKLIŲ SKAIČIUS IR KOMBINATORIKA
% ==========================================
\section{Kiek daliklių turi skaičius? (Labas, Kombinatorika!)}

Dažnas olimpiadų klausimas: kiek daliklių turi skaičius $1 000 000$? Bandyti juos išrašyti ranka būtų savižudybė. Bet čia į pagalbą ateina mūsų sena draugė – Kombinatorika.

Tarkime, turime skaičių $N$, kurio DNR yra $N = p_1^{a_1} \cdot p_2^{a_2}$.
Jei kažkoks skaičius $d$ yra mūsų $N$ daliklis, vadinasi $d$ yra pastatytas \textbf{tik} iš tų pačių „atomų“ $p_1$ ir $p_2$, ir jų kiekis negali viršyti to, kiek jų turi $N$.

Taigi, formuodami daliklį, mes \textbf{kombinatorikos daugybos taisykle} darome pasirinkimus:
\begin{itemize}
    \item Kiek turime variantų pasirinkti pirminį skaičių $p_1$? Galime jo neimti visai ($0$ vnt.), imti $1$ vnt., $2$ vnt. ir t.t. iki pat $a_1$. Iš viso: \textbf{$(a_1 + 1)$} pasirinkimų.
    \item Kiek turime variantų pasirinkti $p_2$? Analogiškai, nuo $0$ iki $a_2$. Iš viso: \textbf{$(a_2 + 1)$} pasirinkimų.
\end{itemize}

Vadinasi, bendras visų įmanomų daliklių skaičius (žymimas graikiška raide $\tau(n)$ - „tau“) yra lygus šių pasirinkimų sandaugai!

\begin{center}
\fbox{
\begin{minipage}{0.9\textwidth}
\textbf{Daliklių skaičiaus formulė:}\\
Jei skaičiaus DNR yra $n = p_1^{a_1} \cdot p_2^{a_2} \cdot \dots \cdot p_k^{a_k}$, tai to skaičiaus daliklių kiekis yra:\\
\[ \tau(n) = (a_1 + 1)(a_2 + 1) \dots (a_k + 1) \]
\textit{Pastaba: Mes dauginame ne pačius pirminius skaičius, o jų \textbf{laipsnius, padidintus vienetu}!}
\end{minipage}
}
\end{center}

\textbf{Pavyzdys:} Kiek daliklių turi skaičius $72$? 
Skaidome: $72 = 8 \cdot 9 = 2^3 \cdot 3^2$.
Daliklių skaičius: $(3+1)(2+1) = 4 \cdot 3 = \mathbf{12}$.

% ==========================================
% TOBULIEJI KVADRATAI IR LYGINUMAS
% ==========================================
\section{Daliklių poros ir stebuklingi kvadratai}

Pabandykime išrašyti skaičiaus $12$ visus daliklius: $1, 2, 3, 4, 6, 12$. 
Pastebite dėsningumą? Dalikliai \textbf{visada vaikšto poromis}! 
Jei išsirinksime bet kokį daliklį $d$, jį padauginę iš „porininko“ ($\frac{12}{d}$) visada gausime pradinį skaičių:
$1 \cdot 12 = 12, \quad 2 \cdot 6 = 12, \quad 3 \cdot 4 = 12$.
Kadangi jie vaikšto poromis, daliklių skaičius paprastai būna \textbf{lyginis}.

Bet ar kada nors gali nutikti taip, kad daliklių skaičius bus \textbf{nelyginis}?
Tai įmanoma tik tuo atveju, jei koks nors daliklis $d$ poroje turės \textbf{pats save}! 
T.y., kai $d \cdot d = N$, arba $N = d^2$. 
O toks skaičius matematikoje vadinamas \textbf{tobuluoju kvadratu} ($1, 4, 9, 16, 25 \dots$).

\textbf{Olimpiadinis faktas:} Natūralusis skaičius turi \textbf{nelyginį} daliklių skaičių tada ir tik tada, kai tas skaičius yra tobulasis kvadratas.

% ==========================================
% PAVYZDŽIAI
% ==========================================
\section{Olimpiadiniai triukai su pirminiais skaičiais}

\begin{example}[Klasikinis algebrinis įrodymas]
Įrodykite, kad jei $p$ yra pirminis skaičius, didesnis už $3$ (t.y. $p \ge 5$), tai reiškinys $p^2 - 1$ \textbf{visada} dalijasi iš $24$.
\end{example}

\begin{solution}
\textbf{Sprendimas:}\\
Šis uždavinys – absoliuti 9-10 klasių klasika. Išskaidykime reiškinį dauginamaisiais (taikydami kvadratų skirtumo formulę):
\[ p^2 - 1 = (p - 1)(p + 1) \]
Pastebėkime, kad $(p-1), p, (p+1)$ yra trys paeiliui einantys sveikieji skaičiai.

\textbf{1 žingsnis (Dalumas iš 8):}
Kadangi $p \ge 5$ yra pirminis, jis privalo būti nelyginis. Vadinasi, jo kaimynai $(p-1)$ ir $(p+1)$ yra \textbf{lyginiai} skaičiai. Mes turime dviejų iš eilės einančių lyginių skaičių sandaugą. Jei paimsite bet kokius du iš eilės einančius lyginius skaičius (pvz., 2 ir 4, arba 10 ir 12), vienas iš jų visada dalijasi iš 2, o kitas – \textbf{būtinai dalijasi iš 4}. Todėl jų sandauga visada dalijasi iš $2 \cdot 4 = 8$. Vadinasi, $8 \mid (p^2 - 1)$.

\textbf{2 žingsnis (Dalumas iš 3):}
Iš bet kokių trijų paeiliui einančių skaičių \textbf{lygiai vienas} visada dalijasi iš 3. Kadangi $p$ yra pirminis (ir $p \ge 5$), pats $p$ iš 3 dalintis negali. Vadinasi, iš 3 privalo dalintis arba $(p-1)$, arba $(p+1)$. Bet kuriuo atveju, jų sandauga dalinsis iš 3. Vadinasi, $3 \mid (p^2 - 1)$.

\textbf{Išvada:}
Kadangi $p^2 - 1$ dalijasi ir iš 8, ir iš 3, o skaičiai 8 ir 3 neturi jokių bendrų daliklių (yra tarpusavyje pirminiai), tai $p^2 - 1$ privalo dalintis iš jų sandaugos: $8 \cdot 3 = \mathbf{24}$. \hfill $\square$
\end{solution}


\begin{example}[Skaičiaus „DNR“ paieška]
Raskite mažiausią natūralųjį skaičių, turintį lygiai $15$ daliklių.
\end{example}

\begin{solution}
\textbf{Sprendimas:}\\
Pasinaudokime daliklių skaičiaus formule $\tau(n) = (a_1 + 1)(a_2 + 1)\dots$
Mums žinoma, kad ši sandauga turi būti lygi $15$. 
Kaip galime gauti skaičių $15$ sudauginę sveikuosius skaičius? Yra tik du būdai:
 $15 \cdot 1$ arba  $3 \cdot 5$

\textit{1 atvejis:} Jei formulėje tėra vienas skliaustas $(a_1 + 1) = 15$, tai laipsnis $a_1 = 14$. 
Skaičiaus „DNR“ atrodo taip: $n = p_1^{14}$.
Kad skaičius būtų kuo mažesnis, imame mažiausią įmanomą pirminį skaičių $p_1 = 2$.
Gauname: $2^{14} = 16384$.

\textit{2 atvejis:} Jei formulėje yra du skliaustai $(a_1 + 1)(a_2 + 1) = 3 \cdot 5$.
Tuomet vienas laipsnis yra $2$, o kitas $4$. 
Skaičiaus „DNR“ atrodo taip: $n = p_1^2 \cdot p_2^4$.
Kad skaičius būtų kuo mažesnis, didesnį laipsnį (ketvirtąjį) priskiriame mažiausiam pirminiam skaičiui ($2$), o mažesnį laipsnį (kvadratą) – sekančiam pirminiam skaičiui ($3$).
Gauname: $n = 3^2 \cdot 2^4 = 9 \cdot 16 = \mathbf{144}$.

Kadangi $144$ yra gerokai mažesnis už $16384$, mažiausias skaičius, turintis lygiai $15$ daliklių, yra $144$. \hfill $\square$
\end{solution}

\vspace{1cm}

% ==========================================
% UŽDAVINIAI
% ==========================================
\newpage
\section{Uždaviniai savarankiškam sprendimui (18 iššūkių)}

\begin{enumerate}
    \renewcommand{\labelenumi}{\textbf{\theenumi.}}
    \item Raskite du pirminius skaičius $p$ ir $q$, jei žinoma, kad jų suma $p + q = 101$.
    \item Raskite pirminį skaičių $p$, jeigu žinoma, kad $p^2 + 11$ taip pat yra pirminis skaičius.
    \item Kiek daliklių turi skaičiai: a) $1000$; \quad b) $2026$; \quad c) $2025$?
    \item Kuris skaičius turi daugiau daliklių: $5^{20} \cdot 3^{10}$ ar $2^{15} \cdot 7^{15}$?
    \item Raskite mažiausią natūralųjį skaičių, turintį lygiai $14$ daliklių.
    \item Įrodykite, kad neegzistuoja toks pirminis skaičius $p$, su kuriuo skaičiai $p, p+2$ ir $p+4$ visi trys kartu būtų pirminiai, išskyrus vienintelį atvejį, kai $p=3$.
    \item Koridoriuje yra $100$ uždarytų spintelių, sunumeruotų nuo $1$ iki $100$. Pirmas mokinys praeidamas atidaro visas spinteles. Antras mokinys praeidamas pakeičia kas antros spintelės būseną (2, 4, 6... - jei buvo atidaryta, uždaro, jei uždaryta, atidaro). Trečias mokinys keičia kas trečios spintelės būseną (3, 6, 9...), ketvirtas – kas ketvirtos ir t.t., kol praeina $100$ mokinių. Kurios spintelės liks atidarytos po visų $100$ mokinių praėjimo? Kodėl?
    \item Įrodykite, kad bet kokiems dviem pirminiams skaičiams $p \ge 5$ ir $q \ge 5$, reiškinys $p^2 - q^2$ visada dalijasi iš $24$.
    \item Skaičius baigiasi vienuolika nulių. Koks mažiausias daliklių skaičius, kurį jis gali turėti?
    \item Keliu nulių baigiasi skaičius, kuris gaunamas sudauginus visus natūraliuosius skaičius nuo $1$ iki $100$ (t.y. $100!$)?
    \item Raskite visus pirminius skaičius $p, q, r$, tenkinančius lygtį: $p \cdot q \cdot r = 5(p + q + r)$.
    \item Jei žinoma, kad $p$ ir $8p^2 + 1$ yra pirminiai skaičiai, raskite $p$. 
    \item Ar egzistuoja toks natūralusis skaičius $n$, kad jo daliklių skaičius $\tau(n)$ būtų lygus pačiam skaičiui $n$? Jei taip, raskite juos visus.
    \item Raskite visus natūraliuosius skaičius $x$ ir $y$, tenkinančius lygtį $x^2 - y^2 = 31$.
    \item Raskite visus natūraliuosius skaičius $x$ ir $y$, tenkinančius lygtį $xy + 2x - 3y = 20$.
    \item Įrodykite, kad skaičius $2^{100} + 5^{100}$ nėra pirminis.
    \item Raskite mažiausią skaičių, kuris pusiau padalintas tampa tobuluoju kvadratu, o padalintas iš 3 – tobuluoju kubu.
    \item \textbf{Iššūkis:} Tegu $d_1, d_2, \dots, d_k$ yra visi skaičiaus $N$ dalikliai (įskaitant 1 ir patį $N$). Įrodykite, kad sandauga $d_1 \cdot d_2 \cdot \dots \cdot d_k = N^{\frac{k}{2}}$.
\end{enumerate}

\end{document}